\documentclass[11pt]{article}

\usepackage[T1]{fontenc}
\usepackage[utf8]{inputenc}
\usepackage{lmodern}
\usepackage[margin=1in]{geometry}
\usepackage{microtype}
\usepackage{mathtools}
\usepackage{amssymb}
\usepackage{amsthm}
\usepackage{booktabs}
\usepackage{longtable}
\usepackage{array}
\usepackage{enumitem}
\usepackage{xcolor}
\usepackage{hyperref}

\definecolor{linkblue}{HTML}{1F4E79}
\hypersetup{
  colorlinks=true,
  linkcolor=linkblue,
  citecolor=linkblue,
  urlcolor=linkblue,
  pdftitle={Actuality Without Exhaustion},
  pdfauthor={Latent-Actual Project},
  pdfsubject={A Lean-checked structural metaphysics of latency and actuality},
  pdfkeywords={metaphysics, actuality, potentiality, process, Lean 4}
}

\setlength{\emergencystretch}{3em}
\setlist{itemsep=0.35em, topsep=0.5em}
\urlstyle{tt}

\newtheorem{formalresult}{Formal result}[section]
\newtheorem{proposition}[formalresult]{Proposition}
\theoremstyle{definition}
\newtheorem{definition}[formalresult]{Definition}
\theoremstyle{remark}
\newtheorem*{qualification}{Qualification}

\newcommand{\lean}[1]{\nolinkurl{#1}}
\newcommand{\Art}{\mathsf{Art}}
\newcommand{\Real}{\mathsf{Real}}
\newcommand{\Res}{\mathsf{Res}}
\newcommand{\Lat}{\mathsf{Lat}}
\newcommand{\Exc}{\mathsf{Exc}}
\newcommand{\Opn}{\mathsf{Open}}
\newcommand{\Und}{\mathsf{Undrawn}}
\newcommand{\Seed}{\mathsf{seed}}
\newcommand{\Offer}{\mathsf{offer}}
\newcommand{\Trace}{\mathsf{trace}}
\newcommand{\Prom}{\mathsf{promoteNext}}
\newcommand{\Index}{\mathsf{index}}
\newcommand{\Active}{\mathsf{Active}}
\newcommand{\After}{\mathsf{after}}
\newcommand{\opp}{\mathsf{opp}}
\newcommand{\pole}{\mathsf{pole}}
\newcommand{\negd}[1]{\overline{#1}}

\title{\textbf{Actuality Without Exhaustion}\\[0.45em]
  \large Opposition, Focus, and Recurrence Without Return}
\author{Latent--Actual Project}
\date{Revision 11 formalization, August 2026}

\begin{document}

\maketitle

\begin{abstract}
This paper develops a Lean-checked structural metaphysics of actuality and
latency. Determinations inhabit unlabelled two-pole oppositions; presentations
distinguish articulation from realization; and continuations settle a target
while carrying a fresh trace of their predecessor. Under an explicit
plenitude postulate, the possible is locally productive and exceeds every
single selection whenever a presentation is supplied, whereas a focus defines
a productive, target-determinate actual sublayer. On histories rooted in one
effective seed, realized, excluded, open, and undrawn content coexist at every
reachable presentation.
An articulation-complete trace-encoding law marks archival precedence within later
presentations as a strict past. A two-beat cadence makes each promotion
followed by reversal, yet this past grows, yielding reversal -- and,
in the wider possible layer, restricted qualitative recurrence -- without
return. An absolutely null countermodel shows that the formalism does not
derive actuality from presentation alone. Its achievement is therefore
conditional and structural: it identifies assumptions under which determinacy,
novelty, non-exhaustion, internal succession, and historical irreversibility
are jointly coherent.
\end{abstract}

\noindent\textbf{Keywords:}
actuality; latency; opposition; process metaphysics; modality; recurrence;
formalization; Lean 4

\tableofcontents
\newpage

\section{The problem}

Metaphysics repeatedly confronts a cluster of tensions. If actuality is
determinate, how can it remain open to more than it presently is? If change
preserves enough structure to count as continuation, how can it be genuinely
novel? If a quality returns, must the world itself return? And if several
continuations are possible, what distinguishes the one that is actual without
reducing possibility to a disguised list of actual facts?

The distinction between potentiality and actuality is ancient, but the present
account is not an exegesis of Aristotle's \emph{Metaphysics} \cite{Aristotle2016}
or of any later system. It isolates a narrower structural question:

\begin{quote}
Can one give a consistent, explicit account in which actuality is locally
determinate and generative, latency remains permanent, qualitative reversal
coexists with historical advance, and the possible exceeds every single
actual orientation?
\end{quote}

The answer formalized here is conditional but positive. The conditions are:

\begin{enumerate}[label=(\roman*)]
  \item a family of binary, unlabelled oppositions;
  \item presentations that distinguish articulation from realization;
  \item continuations that settle a target, preserve old articulation, and
        carry a fresh trace;
  \item an optional plenitude principle for lawful continuations;
  \item a total focus that selects the target of actual transitions; and
  \item one realized seed at the root of the history under consideration.
\end{enumerate}

From these premises the formal development proves local productivity,
target determinacy, permanent latent excess, a strict archive order, finite
no-return, cadence parity, and the swirl: every promotion is followed by
realization of the promoted determination's contrary while the archive moves
on.

This result has a deliberately modest modal force. It is not a proof that the
world has the proposed structure. It is not a derivation of being from
non-being. It is a representation theorem: if the primitives and transition
laws have the stated form, then a distinctive package of metaphysical
consequences follows. The Lean formalization makes the conditional character
useful rather than evasive, because it records exactly which conclusions use
which premises and supplies countermodels where stronger conclusions fail.

\section{Formal method and interpretive discipline}

\subsection{Why formalize a metaphysical argument?}

Formalization cannot decide whether a metaphysical primitive is adequate to
experience or nature. It can, however, perform three kinds of conceptual work.
First, it prevents equivocation between nearby notions. Here
\emph{realized}, \emph{articulated}, \emph{open}, and \emph{possible} have
different types and proof conditions. Second, it exposes hidden premises.
No-return, for example, is not obtained from bare opposition; it depends on
fresh traces, their carriage, and monotone articulation. Third, it separates
implication from interpretation. Lean verifies the implication. The
metaphysical reading remains an argument whose scope can be assessed
independently.

The companion source is \lean{LatentActual.lean}, revision 11. It imports only
Lean's core \lean{Init} module and is checked with Lean 4.24.0
\cite{deMouraUllrich2021}. The concrete
model and principal regression routes contain no user-declared
\lean{axiom}, \lean{sorry}, \lean{admit}, or \lean{native_decide}; executable
\texttt{\#print axioms} guards report only \lean{propext} and
\lean{Quot.sound} for
the checked declarations.

\subsection{Three levels}

The theory is easiest to understand when three levels are kept separate.

\begin{description}
  \item[Determinations.]
  These are poles within oppositions. Contrary determination belongs to the
  same opposition.

  \item[Presentations.]
  These assign articulation and realization status to determinations and carry
  a seed, offer, trace, and cadence bit.

  \item[Transition layers.]
  These relate a source presentation, successor presentation, and settled
  determination. The possible layer records lawful continuations; the actual
  layer is the part of that relation selected by a focus.
\end{description}

This separation blocks two tempting slides. First, the contrary operation on a
determination is not propositional negation. Second, there are two technical
senses near the ordinary word \emph{actual}: \(\Real_p(d)\) is a predicate
saying that content \(d\) is realized in presentation \(p\), while an actual
transition is an edge in a focused subrelation. A realized content is not
itself a transition, and an actual transition is defined through more than the
realization predicate.

\subsection{Conditionality as a result}

The account labels its strong principles instead of hiding them. The open
offer, fresh trace, and articulation-complete trace-encoding law are fields of the
presentation structure. Trace carriage and cadence are fields of
continuation. Plenitude is an optional system-level capability. Rooted seed
effectiveness is a local historical premise. Thus the formal result does not
say that opposition alone produces time, novelty, or being. It says which
package of assumptions suffices, how the pieces interact, and where they do
not.

\section{Opposition before orientation}

\subsection{The field}

A \lean{Field} consists of a type \(O\) of oppositions and, for each
\(o\in O\), a type \(P(o)\) of poles. It provides a pole exchange
\[
  \mathsf{flip}_o : P(o)\longrightarrow P(o)
\]
with no fixed point, together with the dichotomy
\[
  \forall p,q\in P(o),\qquad p=q\ \lor\ p=\mathsf{flip}_o(q).
\]
Equality of oppositions is decidable. A determination is a dependent pair
\[
  d=(\opp(d),\pole(d)),\qquad \pole(d)\in P(\opp(d)).
\]
Its contrary is
\[
  \negd d=(\opp(d),\mathsf{flip}(\pole(d))).
\]

The two-pole dichotomy and fixed-point freedom imply, rather than assume, that
flip is involutive:
\[
  \negd{\negd d}=d,\qquad \negd d\neq d.
\]
These are \lean{Field.neg_involutive} and \lean{Field.neg_free}. If two
determinations share an opposition, then one is either the other or its
contrary (\lean{Field.eq_or_neg}).

\begin{qualification}
A bare \lean{Field} does not require every \(P(o)\) to be inhabited. Strictly
speaking, an inhabited fibre is a free transitive two-element action, while an
empty fibre satisfies the universally quantified field laws vacuously. A
total selector in a full system supplies a pole at every opposition it is
asked to select. Claims that all field fibres are two-pole torsors should
therefore be read with this inhabitation qualification.
\end{qualification}

\subsection{Unlabelled opposition and orientation}

The field is prior to a global labelling of its poles. Nothing in its
definition calls one pole positive, true, or privileged. A
\lean{Polarisation} is additional structure: a dependent section choosing a
pole at every opposition. Exchanging every choice gives the opposite
polarisation, and exchanging twice restores the first.

This order of explanation matters. Difference is primitive at the field level;
orientation is supplied later. The theory can therefore distinguish the fact
that a determinable has two poles from the act of taking one pole as the
selected outcome in a presentation.

\section{Presentation and stratified latency}

\subsection{Articulation, realization, and resolution}

For a presentation \(p\) and determination \(d\), write
\[
  \Art_p(d) \quad\text{when \(p\) articulates \(d\),}
\]
\[
  \Real_p(d) \quad\text{when \(p\) realizes \(d\).}
\]
The presentation laws require:

\begin{enumerate}[label=(P\arabic*)]
  \item articulation is closed under contrary:
        \(\Art_p(d)\Rightarrow\Art_p(\negd d)\);
  \item realization implies articulation:
        \(\Real_p(d)\Rightarrow\Art_p(d)\);
  \item realization is exclusive:
        \(\Real_p(d)\Rightarrow\neg\Real_p(\negd d)\);
  \item each presentation has a fresh trace not articulated there;
  \item if \(q\) articulates \(\Trace(p)\), it articulates everything that
        \(p\) articulates; and
  \item each presentation offers an articulated but unresolved opposition.
\end{enumerate}

Resolution is not a further primitive:
\[
  \Res_p(d)\;:\!\!\Longleftrightarrow\;
  \Real_p(d)\lor\Real_p(\negd d).
\]
Because contrary determinations share an opposition,
\(\Res_p(d)\) and \(\Res_p(\negd d)\) are equivalent in effect.

\begin{formalresult}[Resolution orients exactly one pole]
If \(\Res_p(d)\), then exactly one of \(\Real_p(d)\) and
\(\Real_p(\negd d)\) holds.
\end{formalresult}

\begin{proof}[Proof sketch]
Resolution supplies one of the two realization witnesses. Realization
exclusivity rules out the other. In the second case, involutivity converts
exclusivity at \(\negd d\) into exclusion of \(d\). This is
\lean{PresentationLayer.resolution_orients}.
\end{proof}

The result is local and typed. It does not assert a global law of excluded
middle for arbitrary propositions. It says that when a particular
opposition is resolved in a presentation, the presentation realizes one and
only one of its poles.

\subsection{Four status notions}

The formal vocabulary distinguishes:
\begin{align*}
  \Lat_p(d) &:\!\!\Longleftrightarrow \neg\Real_p(d),\\
  \Exc_p(d) &:\!\!\Longleftrightarrow \Real_p(\negd d),\\
  \Opn_p(d) &:\!\!\Longleftrightarrow
     \Art_p(d)\land\neg\Res_p(d),\\
  \Und_p(d) &:\!\!\Longleftrightarrow \neg\Art_p(d).
\end{align*}

\begin{center}
\small
\begin{tabular}{@{}p{0.18\linewidth}p{0.23\linewidth}p{0.23\linewidth}p{0.25\linewidth}@{}}
\toprule
Status & Articulated? & Realized? & Reading \\
\midrule
Realized & yes & \(d\) is realized & effective determination \\
Excluded & yes & \(\negd d\) is realized & latent through realized contrary \\
Open & yes & neither pole is realized & drawn but unsettled \\
Undrawn & no & \(d\) cannot be realized & not articulated there \\
\bottomrule
\end{tabular}
\end{center}

Excluded, open, and undrawn determinations are all latent, as proved by
\lean{excluded_latent}, \lean{open_latent}, and \lean{undrawn_latent}. They are
not treated as interchangeable. In particular, openness is not ignorance about
which pole is secretly realized; it is the non-resolution of an articulated
opposition.

Nor does the development prove that these labels form a constructively
exhaustive partition of every determination at every presentation. Its
positive claim is more precise: along a rooted actual history there are
simultaneous witnesses of realized, excluded, open, and undrawn content.

\subsection{Offer and trace}

Each presentation carries an \emph{offer}, an opposition with some articulated
pole for which neither pole is realized. Hence every presentation contains
open content. Each presentation also carries a \emph{trace} that is not
articulated at that presentation. Hence every presentation contains undrawn
content.

\begin{formalresult}[Non-exhaustion at presentation strength]
Every presentation has an open determination and an undrawn determination.
Consequently, the thesis that every articulated determination is resolved is
inconsistent with the presentation laws.
\end{formalresult}

\begin{proof}[Proof sketch]
The open witness is extracted from \lean{offer_open}; the undrawn witness is
\(\Trace(p)\) by \lean{trace_fresh}. If every articulated determination were
resolved, the offered witness would be both open and resolved. See
\lean{openNonExhaustive_of_offer},
\lean{articulativeNonExhaustive}, and
\lean{fullResolution_inconsistent_of_offer}.
\end{proof}

This is the first sense in which actuality does not exhaust latency. It is
structural, not merely temporal: even before a history is selected, every
admissible presentation contains an unresolved articulation and an undrawn
trace.

\section{Continuation and the archive}

\subsection{What a continuation does}

A continuation \(p\xrightarrow{d}q\), formalized by
\lean{PresentationLayer.Continues}, settles target \(d\) and constrains how
\(q\) may differ from \(p\). Its clauses say:

\begin{enumerate}[label=(C\arabic*)]
  \item \(\Real_q(d)\);
  \item away from \(d\)'s opposition, previously articulated realization
        status is preserved;
  \item prior articulation persists from \(p\) to \(q\);
  \item \(q\) realizes the exact determination \(\Trace(p)\);
  \item genuinely new articulation lies at the target or trace opposition;
  \item away from the target, a newly realized determination is exactly
        \(\Trace(p)\);
  \item \(\Seed(q)=d\); and
  \item \(\Prom(q)=\mathsf{true}\) iff \(d\) was not open at \(p\).
\end{enumerate}

The target clause makes the act effective. The off-target clauses express
locality: an act cannot silently rewrite unrelated, already drawn content.
The articulation clause supplies continuity. The trace clauses produce
historical difference. Seed-writing makes the result of an act available to
the next act. Cadence records whether the next act should promote open content
or turn settled content.

Several apparently stronger properties are derived. Resolution persists
through a continuation; a refreshed offer cannot occupy the just-settled
opposition; and two successors from one source that settle the same target
agree on the realization status of every determination already articulated at
the source. These are \lean{Continues.resolution},
\lean{Continues.offer_avoid}, and \lean{Continues.agree_on_drawn}.

\subsection{A strict archive}

Define
\[
  p\prec q
  \quad:\!\!\Longleftrightarrow\quad
  \bigl(\forall e,\ \Art_p(e)\Rightarrow\Art_q(e)\bigr)
  \land \Art_q(\Trace(p)).
\]
This is \lean{Archive.Precedes}. The first conjunct says that the archive of
articulation is monotone. The second says that \(q\) contains the trace that was
fresh at \(p\).

\begin{formalresult}[Archive advance and finite no-return]
The relation \(\prec\) is irreflexive and transitive. Every active transition
advances it. Therefore no non-empty finite chain of active transitions returns
to its source presentation.
\end{formalresult}

\begin{proof}[Proof sketch]
Irreflexivity follows because \(\Trace(p)\) is not articulated at \(p\).
Transitivity follows from preservation of articulation. Every active edge
inherits the continuation's articulation and trace-carriage clauses, hence
advances \(\prec\). A positive cycle would imply \(p\prec p\), contradicting
irreflexivity. These are \lean{Archive.precedes_irrefl},
\lean{Archive.precedes_trans}, \lean{Archive.active_precedes}, and
\lean{Archive.no_cycle}.
\end{proof}

The proof reveals the price of irreversibility. The arrow is not obtained from
a symmetric state space alone. It is grounded in a principle of historical
individuation: each source has content fresh to it, the successor carries that
content, and articulation is not deleted. This is a substantive premise, not
a theorem of binary opposition.

\subsection{An internally marked past}

Archive precedence is initially a relation used by the theorist to compare two
presentations. The trace-encoding law permits the same order to be recovered
from what a later presentation articulates. Define
\[
  \mathsf{InPast}(q,p)
  \quad:\!\!\Longleftrightarrow\quad
  \Art_q(\Trace(p)).
\]
In words, \(p\) is in the past marked at \(q\) when \(q\) draws \(p\)'s trace.
Trace encoding says that such a marker is complete enough to carry every
articulation of \(p\), not merely a bare name. It does not encode \(p\)'s
realization profile, seed, offer, cadence, or full state.

\begin{formalresult}[Internal past equals archive precedence]
For all presentations \(p,q\),
\[
  \mathsf{InPast}(q,p)\quad\Longleftrightarrow\quad p\prec q.
\]
Consequently, \(\mathsf{InPast}\), read with the later presentation first, is
irreflexive, transitive, and asymmetric.
\end{formalresult}

\begin{proof}[Proof sketch]
From left to right, trace encoding supplies the monotonic-articulation conjunct
of \(p\prec q\), while the definition of \(\mathsf{InPast}\) supplies its trace
conjunct. The reverse implication is immediate from precedence. Strict-order
properties transfer from the archive, with trace freshness also proving
irreflexivity directly. See \lean{Succession.inPast_iff_precedes},
\lean{Succession.inPast_irrefl}, \lean{Succession.inPast_trans}, and
\lean{Succession.inPast_asymm}.
\end{proof}

Three strengths must be distinguished. \(\mathsf{InPast}(q,p)\) requires only
articulation of \(\Trace(p)\). \lean{ObserverLayer.Registers q p} requires the
trace opposition to be resolved, and
\lean{ObserverLayer.EffectivelyRegisters q p} requires realization of the exact
trace pole. An immediate successor effectively registers its source, and a
later endpoint on a positive chain registers the chain source; arbitrary marked
past membership alone implies neither stronger condition.

An active edge \(p\xrightarrow{d}q\) preserves every past marker already drawn
at \(p\) and adds \(p\) itself to \(q\)'s past, while \(p\) is not in its own
past. The theorem does not say that \(p\) is the only newly added past member.
Along a positive chain from \(p\) to \(q\), \(q\) resolves \(p\)'s trace,
inherits the past marked at \(p\), and still has an open offer
(\lean{Succession.past_settled_future_open}).

The development calls the following package \emph{lived succession}. If a
determination \(o\) is realized at \(p\), and an act targets an opposition other
than \(o\)'s, then \(o\) remains realized at \(q\). At \(q\), the earlier
marked past is intact, \(p\) is newly included, \(q\) is not in its own
past, \(p\)'s trace is resolved, and an opposition remains open:
\[
\begin{split}
  \Real_q(o)\;&\land
  \bigl(\forall r,\ \mathsf{InPast}(p,r)\Rightarrow
                       \mathsf{InPast}(q,r)\bigr)\\
  &\land\mathsf{InPast}(q,p)
  \land\neg\mathsf{InPast}(q,q)
  \land\Res_q(\Trace(p))
  \land\exists f\,\Opn_q(f).
\end{split}
\]
These are \lean{Succession.observer_persists} and
\lean{Succession.lived_succession}. Here \emph{observer} means only an
effective determination persisting off target; no appearance relation or
phenomenology is used.

The result adds a minimal registered-past/open-content contrast. Past marking is
strict, the successor is not in its own past, and open content remains at it.
No formal present predicate or future presentation is supplied. The direction
is derived given trace freshness, articulation-complete trace encoding, articulation
preservation, and trace carriage. It is not derived from temporally neutral
opposition, and the partial order is not proved linear.

\subsection{Recurrence without return}

Suppose \(p\xrightarrow{\negd d}q\xrightarrow{d}r\) are possible acts and
\(\Real_p(d)\). Then the two acts restore the realization status, at \(r\), of
every determination articulated at \(p\):
\[
  \forall e,\quad
  \Art_p(e)\Rightarrow
  \bigl(\Real_p(e)\leftrightarrow\Real_r(e)\bigr).
\]
At the same time, \(p\prec r\). This conjunction is
\lean{Archive.recurrence_without_return}.

The recurrence is deliberately restricted. It is status-equivalence on the
starting presentation's drawn content, not equality of presentations and not
even unrestricted agreement on all determinations. The trace makes the later
presentation historically richer. A pattern can recur without the event
being numerically the same event.

\section{Possibility, plenitude, and excess over selection}

\subsection{Generative layers}

A generative layer \(G\) supplies a ternary relation
\[
  G.\Active(p,q,d),
\]
and requires every active triple to satisfy the continuation laws. It is
\emph{productive} when every presentation has some active successor:
\[
  \forall p\ \exists q,d,\quad G.\Active(p,q,d).
\]
This is one-step seriality. By itself it neither constructs an infinite
history nor selects a coherent successor at every stage.

Presentation structure does not imply productivity. The development defines
an inert generative layer with no active transitions. This counterexample,
\lean{Possibility.productivity_is_not_presentation_level}, motivates an
explicit generative premise.

\subsection{The strength of plenitude}

The optional plenitude capability combines \lean{Possibility.Rich} and
\lean{Possibility.Maximal}.

Richness provides:
\begin{enumerate}[label=(R\arabic*)]
  \item a continuation settling every open determination; and
  \item a continuation turning the seed at every presentation.
\end{enumerate}

Let \(\mathsf{Extends}(G,G')\) mean that every \(G\)-active transition is
\(G'\)-active. Maximality says that if \(G'\) extends \(G\), then \(G\) also
extends \(G'\). Since every generative layer is contained in the full
continuation relation, the theorem
\lean{Possibility.maximal_iff_extends_full} shows that maximality is
extensionally fullness:
\[
  G\text{ is maximal}
  \quad\Longleftrightarrow\quad
  \mathsf{full}\subseteq G.
\]

\begin{qualification}
Plenitude is therefore a strong name for a strong postulate. It does not
follow from opposition or presentation, and it is not merely the absence of a
larger chosen possibility relation. It makes every lawful continuation active.
\end{qualification}

Rich seed-turn continuations transferred by maximality yield productivity
(\lean{Possibility.productive_of_plenitude}). The proof is local: for any
source, turn its seed.

\subsection{Why possibility exceeds every selector}

A selector chooses one determination at each presentation and opposition:
\[
  S(p,o)\in P(o).
\]
Its lifted choice on a determination depends only on that determination's
opposition. Thus
\[
  S_p(d)=S_p(\negd d).
\]
A selector \emph{respects} a generative layer when every active target agrees
with its source selection.

\begin{formalresult}[Plenitude excludes universal selection]
If the presentation layer is rich, \(G\) is maximal, and at least one
presentation is supplied, then no selector is respected by every
\(G\)-active transition.
\end{formalresult}

\begin{proof}[Proof sketch]
At the supplied presentation, choose a pole \(d\) at the open offered opposition.
Openness is symmetric across the two poles. Richness supplies one continuation
settling \(d\) and another settling \(\negd d\); maximality transfers both into
\(G\). If one selector respected both, it would have to choose both \(d\) and
\(\negd d\) at the same presentation and opposition. Coherence says those
selector queries are equal, while fixed-point freedom says
\(d\neq\negd d\). See \lean{Possibility.plenitude_excludes_selection}.
\end{proof}

This is modal excess, not contradiction. A rooted system supplies the needed
presentation witness. The two poles occur as targets of
alternative continuations from one source. They are not co-realized in one
presentation. The possible outruns any one orientation because possibility
branches across contrary lawful acts.

\section{Focus and determinate actuality}

\subsection{From a selector to a focus}

A \lean{Focus} contains a total selector with one constraint: at the seed's
opposition, it selects the seed's contrary. A seeded global polarisation is one
way to build such a focus, but the dynamic core does not require the selector
to be uniform across presentations.

The cadence bit chooses the opposition under consideration:
\[
  \mathsf{focusOpp}(p)=
  \begin{cases}
    \Offer(p), & \Prom(p)=\mathsf{true},\\
    \opp(\Seed(p)), & \Prom(p)=\mathsf{false}.
  \end{cases}
\]
The target index is the focus's selection at this opposition:
\[
  \Index_\Phi(p)=\Phi.\mathsf{select}
    \bigl(p,\mathsf{focusOpp}(p)\bigr).
\]
Consequently, in turn cadence
\[
  \Index_\Phi(p)=\negd{\Seed(p)},
\]
while in promotion cadence its opposition is the open offer.

Given a possible layer \(G\), actuality is the focused subrelation
\[
  (\Phi.\mathsf{actual}(G)).\Active(p,q,d)
  \quad:\!\!\Longleftrightarrow\quad
  G.\Active(p,q,d)\land d=\Index_\Phi(p).
\]
Actuality does not abolish modal alternatives by declaring them incoherent.
It restricts a broader relation by a state-relative orientation.

\subsection{What determinacy amounts to}

Actuality respects its focus by construction
(\lean{Focus.actual_respected}). If two actual edges leave the same source,
their targets are equal. Moreover, their successors agree on the realization
status of everything articulated at the source.

\begin{formalresult}[Limited actual determinacy]
For actual edges \(p\xrightarrow{d}q\) and
\(p\xrightarrow{d'}q'\),
\[
  d=d'
\]
and
\[
  \forall e,\quad \Art_p(e)\Rightarrow
  \bigl(\Real_q(e)\leftrightarrow\Real_{q'}(e)\bigr).
\]
\end{formalresult}

\begin{proof}[Proof sketch]
Both active targets equal the single index \(\Index_\Phi(p)\). With the targets
identified, continuation agreement on previously drawn content applies. See
\lean{Focus.actual_index_unique},
\lean{Focus.actual_target_determinate}, and the bundled theorem
\lean{LatentActualSystem.actual_determinate}.
\end{proof}

The theorem does not prove \(q=q'\). New open articulations and other
presentation structure may differ within the locality constraints. Actuality
is target-determinate without being a unique-future transition system.

Under plenitude the chosen target always has a lawful possible continuation:
when promoting, selection preserves openness at the offer; when turning,
richness supplies the contrary seed continuation. Thus focused actuality is
productive (\lean{Focus.actual_productive}). This dependence matters.
Selection alone chooses a target but does not make a matching successor exist.

\subsection{Cadence}

If an act settles open content, the successor enters turn cadence. If an act
turns a realized contrary, the successor enters promotion cadence. Along a
history whose seed is effective at turn positions, every actual edge flips the
Boolean cadence:
\[
  \Prom(q)=\mathsf{not}(\Prom(p)).
\]

For a finite actual chain of length \(n\),
\[
  \Prom(q)=\After(n,\Prom(p)),
\]
where \(\After\) negates the bit once per act. Hence even-length runs restore
the initial cadence and odd-length runs reverse it. These results are
\lean{Focus.actual_chain_cadence} and the system-level theorems
\lean{actual_chain_cadence_even} and \lean{actual_chain_cadence_odd}.

The Boolean is a minimal control structure, not yet a phenomenology of time.
Its philosophical use is to make a structural alternation exact: openness is
promoted, settled determination is turned, and the cycle repeats at the level
of regime.

\subsection{The swirl}

\begin{formalresult}[Promotion, reversal, and archive advance]
Suppose
\[
  p\xrightarrow{d}_{\mathrm{actual}}q
  \xrightarrow{e}_{\mathrm{actual}}r
\]
and \(p\) is in promotion cadence. Then
\[
  e=\negd d
  \qquad\text{and}\qquad
  p\prec r.
\]
\end{formalresult}

\begin{proof}[Proof sketch]
The first edge selects a pole of the open offer. Continuation cadence therefore
puts \(q\) in turn mode and writes \(\Seed(q)=d\). Focus at a turn selects
\(\negd{\Seed(q)}=\negd d\), so the second target is the contrary of the first.
Each active edge advances the archive, and transitivity gives \(p\prec r\).
See \lean{LatentActualSystem.actual_promote_then_reverse}; the parity-indexed
run version is \lean{LatentActualSystem.Rooted.actual_chain_swirl}.
\end{proof}

The theorem deserves both emphasis and restraint. It proves an actual
one-step reversal after every promotion. It does not prove that one opposition
oscillates forever: after the turn, the next promotion concerns the refreshed
offer, which is away from the settled seed opposition. Nor does the theorem by
itself prove full status recurrence. The stronger
\lean{recurrence_without_return} result belongs to a suitable two-edge
\emph{possible} sequence and is restricted to previously articulated content.

\section{Rooted history and permanent latency}

\subsection{One effective seed}

A rooted system chooses an initial presentation \(p_0\) and assumes
\[
  \Real_{p_0}(\Seed(p_0)).
\]
Every active successor realizes its target and writes that target as its seed.
The initial fact therefore propagates along every finite actual chain from the
root:
\[
  \mathsf{Reachable}(p)\Rightarrow\Real_p(\Seed(p)).
\]
This is \lean{LatentActualSystem.reachable_seed_realised}.

The rooted design is premise-reducing. The system does not assume that every
presentation, including unreachable ones, realizes its seed. It assumes one
effective root and derives the invariant only where actual history reaches.

\subsection{The strongest coexistence result}

\begin{formalresult}[Permanent latency on rooted actual history]
At every reachable presentation \(p\), there exist determinations witnessing:
\[
  \exists d\,\Real_p(d),\qquad
  \exists e\,\Exc_p(e),\qquad
  \exists f\,\Opn_p(f),\qquad
  \exists g\,\Und_p(g).
\]
\end{formalresult}

\begin{proof}[Proof sketch]
The realized witness is \(\Seed(p)\), whose effectiveness propagates from the
root. Its contrary is excluded. The offered opposition supplies an open
witness. The fresh trace supplies an undrawn witness. This is exactly
\lean{LatentActualSystem.latency_is_permanent}.
\end{proof}

The conclusion is not that actuality and possibility are vague mixtures.
Realization remains exclusive and actual targets remain determinate. Rather,
determinacy is local: some content is effective, while contrary, unresolved,
and unarticulated content persists around it. Actuality is bounded without
being deficient. It is a determination within a field that is not totally
resolved.

\subsection{The void and the limit of the anti-nullity argument}

The development also constructs a \lean{Void.layer}: a coherent presentation
layer with an open offer and a fresh trace but no realized determination at
all. It proves
\[
  \exists F,P,p,\quad \forall d,\ \neg\Real_p(d).
\]
Thus absolute nullity is compatible with presentation structure. This is
\lean{nullity_failure_not_derivable}.

The countermodel blocks an inflated reading of the rooted theorem. The account
does not derive something from nothing. It shows that once one seed is
effective, the continuation laws preserve non-nullity along its history while
never exhausting latency. The origin of effectiveness remains a premise.

\section{The concrete Stage model}

Abstract implication is not enough to show that all the assumptions can hold
together. The \lean{Stage} namespace supplies a concrete model.

Its opposition indices have three forms:
\[
  \mathsf{base},\qquad \mathsf{opp}(k),\qquad \mathsf{tr}(m),
\]
with Boolean poles. A stage records a base pole, a partial map of settled
numbered oppositions, a clock, an independent history counter, its seed, a
proof that the seed is effective, and the cadence bit. The current numbered
opposition at the clock is the open offer. The current history index is the
fresh trace.

The continuation \(\mathsf{cont}(s,k,a)\) settles pole \(a\) at numbered
opposition \(k\), increments the history, advances the clock far enough to
refresh the offer, writes the new seed, and computes the next cadence from
whether the target was already resolved. The proof
\lean{Stage.cont_continues} verifies every locality, preservation, trace,
seed, and cadence clause.

The model supplies richness, the full generative layer, a seeded focus,
plenitude, and a rooted initial stage. It proves:

\begin{itemize}
  \item possible and actual productivity;
  \item incompatibility of plenitudinous possibility with one respected
        selector;
  \item an initial actual seed turn;
  \item a following actual promotion of a genuinely different opposition;
  \item exact trace locality for new articulation and off-target realization;
  \item a concrete witness in which an off-target base determination persists
        while the predecessor enters the successor's non-reflexive past; and
  \item a possible two-step status recurrence that still advances the archive.
\end{itemize}

The separate history counter is important. Settling a distant numbered
opposition does not silently make a block of unrelated trace tokens effective.
Historical novelty is carried one predecessor at a time.

The Stage model is a satisfiability witness, not an empirical cosmology. Its
natural numbers and Booleans show that the abstract clauses are jointly
inhabited and computationally manageable. They do not establish that physical
or lived reality is discrete, binary, or clocked in this way.

\section{Metaphysical interpretation}

\subsection{A process ontology with local determination}

The formal picture is processual in a precise sense: presentations are ordered
by acts that preserve articulation, settle a target, and generate an
irreducible historical trace. A presentation is not merely a timeless set of
truths, because its place in the archive depends on what predecessor trace it
carries. Yet the process is not flux without identity. Off-target status
preservation and monotone articulation provide continuity.

This combination resembles a recurring aspiration of process metaphysics:
becoming should produce determinate outcomes without reducing novelty to the
rearrangement of a closed inventory. Whitehead's insistence that actual
entities are acts of becoming rather than inert substances is a relevant
comparison \cite{Whitehead1978}. The present formalization is much thinner. It
does not model prehension, causal efficacy, or a theory of relations. Its
contribution is the smaller structural claim that local settlement, preserved
content, and fresh historical difference can coexist consistently.

\subsection{History marked within presentations}

The succession layer strengthens the archive's metaphysical role. Without
trace encoding, a theorist can prove that \(p\prec q\), but the occurrence of
\(\Trace(p)\) within \(q\) need not by itself entail all of \(p\)'s drawn
content. With complete encoding, the trace functions as an
articulation-complete marker: articulating it is equivalent to standing later
in the archive. Historical order is therefore marked in the articulation
profiles of later presentations, rather than available only to a metalanguage
comparing states. This is extensional marking, not a decoding operation or
knowledge predicate.

The theorem named \lean{lived_succession} captures three structural contrasts.
There is preservation, because an untouched effective determination and the
previously marked past persist. There is non-coincidence, because the
successor includes its source but not itself in its marked past. There is
continuing openness, because an opposition remains articulated and unresolved.
This is a useful minimal pattern of realized content amid marked history and
unresolved content; it is not a formal present or a guaranteed future event.

It is not yet lived time in a phenomenological sense. The effective
determination called an observer has no perspective or experience, the strict
past need not be linearly ordered, and the model has no duration, anticipation,
or affect. The result is better understood as a formal precondition for lived
succession: a presentation-relative past marker plus asymmetric advance,
off-target persistence, and openness. A future account of strands could add
the linearity appropriate to a single history without imposing total order on
the whole archive.

\subsection{Latency is stratified, not a second realm}

Latency here is indexed to a presentation. There is no separate warehouse of
possible objects. A determination is latent at \(p\) exactly when it is not
realized there, but the reason for non-realization may differ:

\begin{itemize}
  \item its contrary is effective;
  \item its opposition is articulated and still open; or
  \item it is not drawn into the presentation.
\end{itemize}

This taxonomy prevents two reductions. Latency is not simply exclusion, since
open and undrawn content need no realized contrary. Nor is it simply
indeterminacy, since excluded content is determined against. The formalism
therefore supports a layered notion of non-actuality without reifying all
possibilities as already actual elsewhere.

The account remains extensional: it does not yet distinguish powers,
dispositions, tendencies, or degrees of readiness within the broad negative
predicate \(\Lat_p(d)\). A richer metaphysics of potentiality would need more
structure than non-realization and lawful continuation.

\subsection{Possibility as lawful transition}

Possibility is not interpreted by a Kripke frame of complete worlds. It is the
availability of lawful successor acts from a presentation. This makes
possibility local and directional. Under plenitude, both poles of an open
opposition can head alternative continuations, so no single selection
represents all possible acts.

The result offers a precise reading of modal excess:
\[
  \text{the possible is not one unactualized total description;}
\]
\[
  \text{it is a branching relation that outruns each single orientation.}
\]
Actuality, by contrast, is not the elimination of the other branches from the
possible relation. It is the focused subrelation. This preserves a clean
distinction between modal breadth and actual determination.

The strength of plenitude limits the interpretation. Because maximality is
extensionally fullness, the theorem characterizes the consequences of a
maximally permissive lawful modality. A weaker modal metaphysics could retain
focus, archive, and rooted latency while losing the universal selector-conflict
result or productivity.

\subsection{Repetition with difference}

The distinction between recurrence and return has affinities with philosophical
accounts that oppose qualitative repetition to the reinstatement of an
identical total state. Bergson's contrast between spatialized succession and
duration \cite{Bergson1911}, and Deleuze's account of repetition as productive
of difference \cite{Deleuze1994}, provide useful points of comparison. No claim
of textual equivalence is intended.

The formal contribution is exact:

\begin{enumerate}
  \item a possible two-act reversal can restore old articulated realization
        status;
  \item actual promotion is followed by contrary realization;
  \item in either relevant two-edge case, fresh trace carriage makes the later
        presentation strictly later in the archive.
\end{enumerate}

Thus recurrence is indexed to a chosen observational profile
(\lean{StatusEqOn}), while return would be equality of presentations. The former
does not imply the latter. The archive encodes a difference that survives
qualitative restoration.

\subsection{Individuation and the seed}

The seed is the target written by an act into its successor. It is both a
record of what has just become effective and the pole available for a future
turn. In this narrow sense, the result of one individuation becomes material
for the next. This invites comparison with Simondon's emphasis on metastability
and the retention of unresolved potentials through individuation
\cite{Simondon2020}.

Again, the comparison marks a structural affinity, not a formalization of
Simondon. The present seed has no energetic magnitude, incompatibility
relation, or transductive domain. What is proved is simpler: settlement writes
a state-relative handle that focus can reverse, while a fresh offer keeps
another opposition open.

\subsection{Internal observation, minimally construed}

An optional observer layer relates a presentation, an observer determination,
and an observed determination. If something appears to an observer at a
presentation, both observer and observed determinations are realized there.
Consequently, contrary contents cannot both appear to the same observer there,
and open content does not appear as realized.

This gives a minimal immanence condition: observation is internal to the
presentation rather than a view from outside the system. It does not yet model
perspective, misrepresentation, attention, consciousness, embodiment, or
phenomenal character. The dynamic core does not depend on the observer layer.

The succession namespace uses \emph{observer} in a still thinner sense: any
realized determination whose opposition an act does not target. Its persistence
is derived from off-target preservation and does not require
\lean{ObserverLayer} or an appearance witness. The two notions should not be
conflated: one constrains appearance, while the other supplies a persisting
effective marker through which the structural succession package can be
stated.

\section{Objections, replies, and limits}

\begin{enumerate}[label=\textbf{O\arabic*.}, leftmargin=*, align=left]
  \item \emph{The conclusions are stipulated in the premises.}

  Some are. Open offers, fresh traces, trace carriage, locality, and cadence
  are structural assumptions. The reply is not that Lean makes them
  self-evident. It is that the formalization turns the complaint into a
  dependency audit: one can see exactly which conclusions require them.
  Several burdens are nevertheless reduced. Involution, resolution
  persistence, offer--seed freshness, refreshed-offer avoidance, cadence
  parity, and no-cycle are theorems rather than constructor fields.

  \item \emph{Binary opposition is metaphysically reductive.}

  The field is a local schema for determinables represented by contrary poles,
  not a claim that every property or relation in reality is binary. Multi-pole,
  continuous, or non-oppositional determinables would require another field
  structure. The present theorem states what follows in the binary case.

  \item \emph{Contrary determination has been confused with logical negation.}

  The types prevent this. \(\negd d\) is a pole exchange within
  \(\opp(d)\). Propositional negation occurs separately in statements such as
  \(\neg\Real_p(d)\). Ontic contrariety and logical denial are related by
  realization exclusivity, not identified.

  \item \emph{Latency is merely a negative predicate.}

  At the broadest level, yes: \(\Lat_p(d)\) is non-realization. But the theory
  refines it into exclusion, openness, and undrawnness with different positive
  structural conditions. What it does not yet supply is an intensional theory
  of powers or propensities. That is a genuine extension point.

  \item \emph{Plenitude is too strong.}

  Maximality is extensionally the full lawful continuation relation, so the
  objection is well founded against any attempt to present plenitude as weak.
  The response is modularity. Plenitude is optional and used for productivity
  and modal excess over selection. No-return, conditional swirl, determinacy,
  and much of the archive theory do not require it.

  \item \emph{Focus makes determinism trivial.}

  Focus does fix the target by definition. The non-trivial content is what
  follows when this filter interacts with openness, seed-writing, continuation
  locality, and cadence. Even then, determinacy is limited: the theory proves a
  unique target and agreement on old drawn status, not a unique successor
  presentation.

  \item \emph{The archive simply stipulates the arrow of time.}

  The underlying ingredients -- fresh trace, articulation-complete trace encoding,
  carriage, and monotone articulation -- are indeed
  historical-individuation premises. The strict order and its internal
  marking are derived from them. The theory therefore explains no-return
  by a precise record principle; it does not derive an arrow from temporally
neutral laws.

  \item \emph{An internally marked past is already memory or knowledge.}

  It is not. The equivalence between \(\mathsf{InPast}(q,p)\) and \(p\prec q\)
  is extensional: \(q\)'s articulation profile contains an
  articulation-complete trace marker. There is no decoding function,
  awareness predicate, or proof that an observer experiences or recalls the
  order. Nor does mere past membership imply that every marked trace is
  resolved.

  \item \emph{The non-nullity argument is circular.}

  Rooted history assumes one realized seed. The void countermodel proves that
  this assumption cannot be removed at presentation strength. The valid claim
  is persistence of actuality from an effective root, not creation of
  actuality from absolute nullity.

  \item \emph{A Boolean cadence is not lived time.}

  Correct. It is the smallest automaton expressing alternation of promotion and
  turn. The separate \lean{lived_succession} theorem adds internal retention,
  asymmetric advance, observer persistence off target, and openness ahead.
  Despite its name, it does not formalize phenomenal duration, variable tempo,
  anticipation, affect, or embodied temporality. Space also remains
  outstanding.

  \item \emph{The swirl overstates oscillation.}

  Read narrowly, it does not. It proves that every promotion has one immediately
  following contrary act, with archive advance. It does not prove endless
  alternation on one opposition. The refreshed offer normally moves the next
  promotion elsewhere.

  \item \emph{The Stage model is artificial.}

  It is artificial in the proper sense of a model construction. Its role is to
  witness joint satisfiability and catch inconsistent axiom packages, not to
  supply empirical evidence or a unique intended interpretation.

  \item \emph{The observer theory does not explain experience.}

  It does not. The optional appearance layer proves internal co-realization and
  some exclusivity consequences. The succession theorem separately proves
  off-target persistence for a realized determination called an observer.
  Neither supplies consciousness, error, salience, or phenomenal character.
\end{enumerate}

\section{Conclusion}

The latent--actual formalization gives a conditional answer to a difficult
structural question. A world of unlabelled oppositions can support locally
oriented realization without global closure. Lawful acts can preserve what is
drawn while making a target and a fresh trace effective. A maximally rich
possible relation, when supplied a presentation, can exceed every single
selector, while focus carves out a productive and target-determinate actuality.
From one effective root, realized, excluded, open, and undrawn content coexist
throughout reachable history. Trace encoding marks archive precedence inside
later articulation profiles. Promotion leads to reversal, but trace carriage
prevents return.

The strongest philosophical lesson is methodological as much as ontological.
Determinacy need not mean exhaustiveness; recurrence need not mean identity;
and conditional metaphysical claims become clearer when their generative,
modal, and historical premises are separately named. The void countermodel
keeps the account honest: presentation does not make itself actual. The Stage
model shows, conversely, that actuality, permanent latency, modal excess,
focused selection, and archival irreversibility can inhabit one coherent
structure.

What remains is substantial. The binary schema must be generalized if it is to
address continuous or plural determination. Productivity must be distinguished
from completed infinite succession. The development now marks a strict past
within presentations and proves a limited persistence package named lived
succession, but metric or continuous time, phenomenal duration, and a formal
future relation remain absent. So do space, causation, embodiment, and
full-blooded observation. These are not small omissions, but the formal core
gives future extensions a disciplined base: new premises can be added, and
their real consequences can be proved rather than suggested.

\appendix

\section{Argument-to-theorem map}

\small
\begin{longtable}{@{}>{\raggedright\arraybackslash}p{0.33\linewidth}
  >{\raggedright\arraybackslash}p{0.61\linewidth}@{}}
\toprule
Lean declaration & Role in the argument \\
\midrule
\endfirsthead
\toprule
Lean declaration & Role in the argument \\
\midrule
\endhead
\lean{Field.neg_involutive} and \lean{Field.neg_free}
  & Contrary exchange is involutive and has no fixed point. \\
\lean{PresentationLayer.resolution_orients}
  & A resolved opposition realizes exactly one pole. \\
\lean{PresentationLayer.openNonExhaustive_of_offer}
  & Every presentation contains an open determination. \\
\lean{PresentationLayer.articulativeNonExhaustive}
  & Every presentation contains an undrawn determination, namely its fresh
    trace. \\
\lean{PresentationLayer.fullResolution_inconsistent_of_offer}
  & Universal resolution of articulated content is inconsistent with the open
    offer. \\
\lean{PresentationLayer.trace_encodes}
  & Articulating a source trace carries every determination articulated at that
    source; it is not a complete state encoding. \\
\lean{Continues.resolution}
  & Resolution persists through lawful continuation. \\
\lean{Continues.newly_articulated_is_trace} and
\lean{Continues.off_target_novel_is_trace}
  & Transition locality confines non-target novelty to the predecessor trace. \\
\lean{Archive.no_cycle}
  & No non-empty finite active chain returns to its identical source. \\
\lean{Succession.inPast_iff_precedes} and
\lean{Succession.inPast_asymm}
  & A presentation's past markers are exactly archive predecessors and form an
    asymmetric relation. \\
\lean{Succession.past_grows_strictly}
  & Every active edge preserves earlier past markers and adds at least its
    source. \\
\lean{Succession.past_settled_future_open}
  & A positive chain resolves its starting trace, preserves its starting past,
    and ends with an open offer. \\
\lean{Succession.observer_persists} and
\lean{LatentActualSystem.lived_succession}
  & Off-target realized content persists through one act, conjoined with strict
    past growth, source-trace settlement, and continuing openness. \\
\lean{Archive.recurrence_without_return}
  & A two-step possible reversal restores old drawn status while the archive
    advances. \\
\lean{Possibility.productive_of_plenitude}
  & Richness plus maximality gives a successor at every source. \\
\lean{Possibility.plenitude_excludes_selection}
  & No one selector agrees with both contrary branches made active by
    plenitude. \\
\lean{Focus.actual_respected}
  & Focused actuality agrees with its selector. \\
\lean{Focus.actual_productive}
  & Plenitude supplies a possible continuation matching every focused index. \\
\lean{LatentActualSystem.actual_determinate}
  & Actual edges from one source share a target and agree on previously drawn
    realization status. \\
\lean{LatentActualSystem.reachable_seed_realised}
  & One effective root seed propagates along actual reachability. \\
\lean{LatentActualSystem.latency_is_permanent}
  & Realized, excluded, open, and undrawn witnesses coexist at every reachable
    stage. \\
\lean{LatentActualSystem.actual_chain_cadence_even} and
\lean{LatentActualSystem.actual_chain_cadence_odd}
  & Actual cadence is determined by finite run parity. \\
\lean{LatentActualSystem.actual_promote_then_reverse}
  & Every promotion is followed by contrary realization and archive advance. \\
\lean{LatentActualSystem.Rooted.actual_chain_swirl}
  & The swirl holds at every promoting parity position in a rooted actual run. \\
\lean{nullity_failure_not_derivable}
  & Presentation structure alone permits absolute nullity. \\
\lean{Stage.actual_productive}, \lean{Stage.actual_novelty},
\lean{Stage.lived_succession_witness}, and
\lean{Stage.oscillation_possible}
  & The concrete model witnesses productivity, novelty, marked succession, and
    possible recurrence without return. \\
\bottomrule
\end{longtable}
\normalsize

\section{Premise ledger}

\begin{center}
\small
\begin{tabular}{@{}>{\raggedright\arraybackslash}p{0.27\linewidth}
  >{\raggedright\arraybackslash}p{0.31\linewidth}
  >{\raggedright\arraybackslash}p{0.33\linewidth}@{}}
\toprule
Layer & Supplied structure & Principal derived consequences \\
\midrule
Field
  & Pole exchange, fixed-point freedom, two-pole dichotomy
  & Involution; equality-or-contrary within a fibre \\
Presentation
  & Articulation closure, realization exclusivity, open offer, fresh trace,
    articulation-complete trace encoding
  & Exact orientation of resolution; open and undrawn non-exhaustion; internally
    marked archive order \\
Continuation
  & Target realization, preservation, trace carriage, locality, seed-writing,
    cadence equation
  & Resolution persistence; offer avoidance; archive advance; local agreement \\
Plenitude
  & Open continuations, seed turns, extensional maximality
  & Possible and actual productivity; no universally respected selector \\
Focus
  & Total state-relative selection; contrary selection at seed opposition
  & Unique actual target; promotion/turn index; actual focus-respect \\
Root
  & One initially realized seed
  & Reachable seed effectiveness; turn cadence; reachable non-nullity and
    permanent latency \\
\bottomrule
\end{tabular}
\end{center}

\begin{thebibliography}{9}

\bibitem{Aristotle2016}
Aristotle.
\newblock \emph{Metaphysics}.
\newblock Translated by C. D. C. Reeve. Hackett Publishing, 2016.

\bibitem{Bergson1911}
Henri Bergson.
\newblock \emph{Creative Evolution}.
\newblock Translated by Arthur Mitchell. Henry Holt, 1911.

\bibitem{Deleuze1994}
Gilles Deleuze.
\newblock \emph{Difference and Repetition}.
\newblock Translated by Paul Patton. Columbia University Press, 1994.

\bibitem{deMouraUllrich2021}
Leonardo de Moura and Sebastian Ullrich.
\newblock The Lean 4 theorem prover and programming language.
\newblock In \emph{Automated Deduction -- CADE 28}, pages 625--635.
Springer, 2021.

\bibitem{Simondon2020}
Gilbert Simondon.
\newblock \emph{Individuation in Light of Notions of Form and Information}.
\newblock Translated by Taylor Adkins. University of Minnesota Press, 2020.

\bibitem{Whitehead1978}
Alfred North Whitehead.
\newblock \emph{Process and Reality: An Essay in Cosmology}.
\newblock Corrected edition, edited by David Ray Griffin and Donald W.
Sherburne. Free Press, 1978.

\end{thebibliography}

\end{document}
